\section{Grundlegender Aufbau}
Auf eine Überschrift folgt immer ein Absatz mit Text, \textbf{nie unmittelbar} eine Überschrift niedrigeren Ranges \cite[S.~197]{Balz2008}.\footnote{Hervorhebungen werden durch Kursivschrift oder halbfette Schrift realisiert. Vermieden werden soll Hervorhebung durch Unterstreichen, größere Schrift oder Großbuchstaben \cite[S.~182]{Balz2008}. Ebenfalls ungünstig ist die farbliche Hervorhebung von Wörtern oder Überschriften.}
Daher muss von der Einleitung eine schlüssige Überleitung zum Hauptteil und den anstehenden Unterkapiteln erfolgen.
Falls mehrere Versuche innerhalb des Praktikums durchgeführt werden, sollte an dieser Stelle kurz darauf eingegangen werden.

Ein Hauptkapitel (Kapitel mit nur einer Zählnummer) beginnt jeweils auf einer neuen Seite.
Innerhalb der Gliederung ist zu beachten, dass Gliederungsebene vier nicht überschritten werden sollte.
Die Unterteilung soll möglichst in der dritten Stufe enden \cite[S. 12]{Hering2009}.

Zum Beispiel soll eine Gliederungsebene "`2.2.1.3"' vermieden werden, also wird die Untergliederung bei Ebene "`2.2.1"' beendet.
Desweiteren muss darauf geachtet werden, dass ein Oberpunkt der Gliederung mindestens zwei Unterpunkte besitzt muss \cite[S. 53]{Karm2011}.
Dennoch sollen die Zählnummern innerhalb einer Hierarchiestufe nicht größer als neun werden \cite[S.~12]{Hering2009}.
Die jeweils übergeordneten Kapitel sollen einen Überblick über die Inhalte geben, die in den Unterkapiteln thematisiert sind.
In den Überblicken sollte jedoch kein zusätzliches Wissen vermittelt werden \cite[S.~196]{Balz2008}.

In diesem Dokument wird ein beispielhaftes Inhaltsverzeichnis für einen Praktikumsbericht aufgezeigt, der u.a. im Praktikum Grundlagen der Elektrotechnik an der HSHL angefertigt werden muss.
Die Inhaltsübersicht und das Inhaltsverzeichnis sollen einen Überblick über die Systematik der Arbeit geben und einen "`ersten Eindruck"' vermitteln.
Man sollte sich bewusst sein, dass das Inhaltsverzeichnis die erste Information ist, die ein Leser von der Arbeit aufnimmt \cite[S.~49]{Karm2011}.
Die Gliederung soll bereits als "`roter Faden"' dienen.
Für den Leser ist die Gliederung der wichtigste Schlüssel zum Erfassen des Inhalts.
Wissenschaftliche Arbeiten erfordern ein hohes Maß an Ordnung und Logik, welche sich auch in der Gliederung widerspiegeln muss \cite[S.~9 ~ff.]{Hering2009}.

\noindent
Üblicherweise gliedert sich der inhaltliche Aufbau in Einleitung, Hauptteil und Schluss.
Die Einleitung beinhaltet die Aufgabenstellung und die Motivation der wissenschaftlichen Arbeit (s. Einleitung, S. \pageref{sec:Einleitung}).
Im Hauptteil werden die theoretischen Grundlagen erläutert und der Lösungswegs sowie die Ergebnisse dargestellt.
Die Ergebnisse werden im Schlussteil in Kurzform wiedergegeben und kritisch diskutiert.
Weitere Untergliederungen dieser drei Hauptpunkte können je nach Art der wissenschaftlichen Arbeit und den entsprechenden Anforderungen umgesetzt werden.
In dieser wissenschaftlichen Arbeit werden die benötigten Gliederungspunkte für einen Praktikumsbericht aufgezeigt.

Ein Praktikumsbericht, der die Beschreibung einen Versuchs beinhaltet, sollte unbedingt folgende Unterpunkte beinhalten:
\begin{enumerate}
\item Theoretische Grundlagen 
\item Aufgabenstellung mit Versuchsbeschreibung
\item Versuchsvorbereitung (benötigter Materialien und Beschreibung des Versuchsaufbaus)
\item Versuchsdurchführung
\item Versuchsauswertung (vgl. \cite{Hering2009}, S. 18)
\end{enumerate}

Als allgemeine Gliederung für ingenieurwissenschaftliche Arbeit (z.\,B Bachelor- oder Masterarbeit) empfiehlt sich die Nachfolgende. 

Titelseite\\
Inhaltsverzeichnis
\begin{enumerate}
	\item Einleitung (Einführung in die Aufgabenstellung)
  \item Theoretische Grundlagen
	\item Lösungsweg ggf. Lösungsalternative
  \item Ergebnis
  \item Zusammenfassung (mit anschl. Diskussion der Ergebnisse) und Ausblick
\end{enumerate}
Quellenverzeichnis\\
Anhang\\
Eidesstattliche Erklärung 
